\section{Integer rows for the finite-order certificates}
\label{f60:subsec-row-specification}

For clarity and reproducibility, the following specifies every row
family used by the finite linear certificates.  A row labelled $\rho$
means $\sum_{j,m}a^\rho_{j,m}t_{j,m}\le b^\rho$; unspecified
coefficients vanish.  All indices satisfy $1\le j\le v$ and
$0\le m\le a-j$.  The symbols in the rows are defined above, and
$r,h,l$ are integers.

\begin{description}
\item[\texttt{count}$(r)$.]
For $2\le r\le v$, use
$a_{j,m}=\mathbf{1}_{j=r}$ and $b=\binom vr$.

\item[\texttt{path}$(r)$.]
For $2\le r\le a$, use
$a_{j,m}=-\binom m{r-j}$ and
$b=\binom ar-\binom{n-r+1}{r}$.

\item[\texttt{private}$(r,h)$.]
For $2\le r\le v$ and $0\le h<a$, set
\[
 a_{j,m}=\mathbf{1}_{j=r}\bigl((r-1)m+a-r+1-rh\bigr)_+
 -(v-r+1)\mathbf{1}_{j=r-1}(m-h)_+,
 \qquad b=0.
\]

\item[\texttt{hall}$(r,l)$.]
Take $0\le l\le a$, $0\le r<\min(v,a-l)$, and
$c=v-r-\max(0,l-\delta)\ge0$.  Set
\[
 a_{j,m}=\bigl((r+1)\mathbf{1}_{j=r+1}-c\mathbf{1}_{j=r}\bigr)
                 \binom ml,\qquad
 b=\begin{cases}c\binom al,&r=0,\\0,&r>0.\end{cases}
\]
The implemented domain with $l=0$ uses $1\le r<v$; the omitted
$(r,l)=(0,0)$ case is the already imposed first-layer equality.

\item[\texttt{tail}$(r,h)$.]
For $2\le r\le v$ and $0\le h\le a-r+1$, set
\[
 a_{j,m}=\mathbf{1}_{j=r}b_{r,h}(m)
 -(v-r+1)\mathbf{1}_{j=r-1}\mathbf{1}_{m\ge h},\qquad b=0.
\]

\item[\texttt{release\_upper}$(r,h)$.]
On the same range, set
\[
 a_{j,m}=s_r\mathbf{1}_{j=r-1}(m-h)_+
            -\mathbf{1}_{j=r}U_{r,h}(m),\qquad b=0.
\]

\item[\texttt{tail\_upper}$(r,h)$.]
On the same range, set
\[
 a_{j,m}=s_r\mathbf{1}_{j=r-1}\mathbf{1}_{m\ge h}
            -\mathbf{1}_{j=r}V_{r,h}(m),\qquad b=0.
\]

\item[\texttt{mean}$(r)$.]
For $2\le r\le v$, set
\[
 a_{j,m}=-m\mathbf{1}_{j=r}-\beta_r\mathbf{1}_{j=2},\qquad
 b=-C_r^0-\beta_r\binom v2.
\]

\item[\texttt{union}$(r)$.]
For $2\le r\le v$, set
\[
 a_{j,m}=(a-m)\left(\mathbf{1}_{j=r}
                -\binom{v-1}{r-1}\mathbf{1}_{j=1}\right),\qquad b=0.
\]

\item[\texttt{edge\_lower}$(r)$.]
For $2\le r\le v$, set
\[
 a_{j,m}=-\mathbf{1}_{j=r}+\binom{v-2}{r-2}\mathbf{1}_{j=2},\qquad
 b=-\binom vr+\binom{v-2}{r-2}\binom v2.
\]

\item[\texttt{edge\_upper}$(r)$.]
For $2\le r\le v$ and $D_0>0$, set
\[
 a_{j,m}=D_0\mathbf{1}_{j=r}-G_r\mathbf{1}_{j=2},\qquad
 b=D_0\binom vr-G_r\binom v2.
\]
For $D_0=0$, use the \texttt{count}$(r)$ row instead.

\item[\texttt{assumption}$(k)$.]
Only in a \texttt{conditional} certificate at the same index $k$, use
\[
 a_{j,m}=B_{m,j}(k),\qquad b=-B_{a,0}(k).
\]
\end{description}

These formulas specify a relaxation containing the counts of every
forest after the prescribed choice of $I$.  They do not assert that
every point in the relaxation is realizable by a forest.

